\section{Введение и основные результаты}\label{sec:intro}

\subsection{Задача}

\emph{Хроматическим числом плоскости} $\chi(\R^2)$ называется наименьшее число
цветов, в которые можно окрасить все точки плоскости так, чтобы никакие две
точки на расстоянии ровно~$1$ не оказались одного цвета. Этот вопрос, известный
как \emph{проблема Нельсона--Хадвигера} (ок.~$1950$), до сих пор не решён: с
$2018$~г. известно лишь $5\le\chi(\R^2)\le7$ \cite{deGrey,Soifer}. Аналогично
определяется $\chi(\R^n)$ для пространства любой размерности.

Естественное обобщение --- запрещать не одно расстояние, а целый \emph{отрезок}:
через $\chi\bigl(\R^n,[1,\ell]\bigr)$ обозначим наименьшее число цветов
раскраски, при которой одноцветные точки не реализуют ни одного расстояния из
отрезка $[1,\ell]$, $\ell\ge1$. При $\ell=1$ отрезок вырождается в точку и мы
возвращаемся к классическому $\chi(\R^n)=\chi\bigl(\R^n,[1,1]\bigr)$.
Систематическое изучение версии с интервалом начато первым автором
в~\cite{Ivanov2011}; в плоскости интервальная версия (в эквивалентной форме
запрета $[1-\varepsilon,1+\varepsilon]$) рассматривалась Экзоо
\cite{Exoo2005}; о современном состоянии плоского случая см.~\cite{ChJW}
и~\cite{Parts2023}. Поскольку запрет более широкого множества
расстояний требует не меньше цветов,
\begin{equation}\label{eq:mono}
\chi(\R^n)=\chi\bigl(\R^n,[1,1]\bigr)\ \le\ \chi\bigl(\R^n,[1,\ell]\bigr)
\qquad(\ell\ge1),
\end{equation}
и любая оценка сверху для интервальной версии автоматически оценивает
классическое хроматическое число. Подробнее об истории задачи см.
обзоры~\cite{BogolubskyRaigorodskii,Raigorodskii2001,Raigorodskii2013,Raigorodskii2014}.

Все известные \emph{верхние} оценки $\chi(\R^n)$ в малых размерностях получены
одним геометрическим приёмом --- \emph{периодической раскраской по разбиению
Вороного решётки} (\S\ref{sec:method}): пространство режется на ячейки Вороного
решётки $\Lambda$, и ячейка узла $v$ получает цвет класса смежности $v+\Gamma$
по подрешётке $\Gamma\subset\Lambda$ индекса $k$; раскраска правильна для
отрезка $[1,\ell]$, если \emph{нормированная ширина}
$d=D(\Gamma)/\diam V_0$ --- отношение наименьшего расстояния между
одноцветными ячейками к диаметру ячейки --- превосходит~$\ell$. Наилучшие
прежние значения в размерностях $4$--$10$ принадлежат Арману, Бондаренко,
Прымаку и Радченко~\cite{ABPR} (далее --- АБПР): $49$, $140$, $343$, $1372$,
$2401$, $17253$ и $3^{10}$.

Верхняя оценка в размерности~$4$ имеет историю: $\le54$ (Радоичич, Тот
\cite{RadoicicToth}, решётка $A_4^*$); оценка $\le49$ заявлена без
доказательства Кулсоном и приведена в \cite{CoulsonPayne2007}; полное
доказательство, вместе с проверкой минимальности индекса~$49$ \emph{для
решётки $D_4$}, дано в~\cite{ABPR}. Нижняя граница в этой размерности прошла
путь $6$ (Райский) $\to7$ \cite{Ivanov2006} $\to9$ \cite{ExooIsmailescu}.
В~\cite{ABPR} (открытый вопрос~1) высказана гипотеза, что улучшить $49$
нельзя:
\begin{quote}\itshape\small
<<We believe that the answer is negative for $n=4,5$, i.e.
$\min\chi_s(\E^4,\Lambda,\Lambda',\ell)=49$~\ldots>> (\cite{ABPR}, Open
question~1),
\end{quote}
где минимум берётся по всем решёткам $\Lambda$, подрешёткам $\Lambda'$ и всем
$\ell>0$. Настоящая работа опровергает эту гипотезу.

\subsection{Результаты}

\begin{keybox}
\begin{theorem}[сводная]\label{thm:summary}
\[
\chi(\R^4)\le43,\qquad \chi(\R^5)\le132,\qquad \chi(\R^7)\le1029,\qquad
\chi(\R^9)\le7203,\qquad \chi(\R^{10})\le45619 .
\]
Точнее, $\chi\bigl(\R^n,[1,\ell]\bigr)\le k$ при всех $\ell\le\ell_0$ для
пар $(n,k,\ell_0)$ из таблицы~\ref{tab:main}.
\end{theorem}
\end{keybox}

\begin{table}[htbp]\centering\small
\caption{Пять доказанных оценок. Столбец $\ell_0$ --- доказанная ширина
запрещённого отрезка $[1,\ell_0]$; запись $<x$ означает, что оценка доказана
при каждом $\ell<x$. Прежние оценки --- из~\cite{ABPR}. Каждая строка
проверяется одним маршрутом, указанным в последнем столбце
(таблица~\ref{tab:manifest} в \S\ref{sec:repro} даёт файлы и команды).}
\label{tab:main}
\begin{tabular}{@{}c r r l l l@{}}
\toprule
$n$ & было & стало & $\ell_0$ & доказательство & где \\
\midrule
$4$  & $49$    & $\mathbf{43}$    & $1{,}00411$      & точный рациональный сертификат & теорема~\ref{thm:main43}\\
$5$  & $140$   & $\mathbf{132}$   & $101/100$        & точный рациональный сертификат & теорема~\ref{thm:r5}\\
$7$  & $1372$  & $\mathbf{1029}$  & $103/100$        & точный рациональный сертификат & теорема~\ref{thm:r7ex}\\
$9$  & $17253$ & $\mathbf{7203}$  & $<1{,}0166127$   & точный рациональный сертификат & теорема~\ref{thm:r9ex}\\
$10$ & $3^{10}$& $\mathbf{45619}$ & $<\sqrt{217/214}$& аналитически (продуктовое правило) & теорема~\ref{thm:dim10}\\
\bottomrule
\end{tabular}
\end{table}

Доказательства устроены по двум схемам.

\emph{Явные конструкции с точными сертификатами} ($n=4,5,7,9$). В каждом случае
предъявляется рациональная матрица Грама $Q$ и целочисленная подрешётка, а
неравенство $d>\ell_0$ сводится к конечному списку неравенств между
рациональными числами, проверенных в точной арифметике дробей: точный
квадрат радиуса покрытия ячейки, доказуемо полный список опасных векторов и
для каждого из них сертификат Каруша--Куна--Таккера расстояния до ячейки.
Единый \emph{протокол} такой проверки сформулирован в \S\ref{sec:method}
(предложение~\ref{prop:protocol}), а четыре конструкции --- в
\S\ref{sec:constr}. Решётка индекса~$43$ \emph{эйзенштейнова}
($\Z[\omega]$-модуль ранга~$2$), решётка индекса~$132$ --- общего положения,
конструкции индексов~$1029$ и~$7203$ --- ламинирования эйзенштейновых
раскрасок $E_6^*/343$ и $E_8/2401$; плавающая точка участвует лишь в
\emph{выборе} кандидатов и активных множеств, но не в проверке.

\emph{Продуктовое исчисление ширин} ($n=10$). Для ортогонального
произведения раскрасок с ширинами $d_i$ условие правильности сворачивается в
одно неравенство $\sum_i1/d_i^2\le1$ (предложение~\ref{prop:product}), причём
блоки могут иметь разный индекс. Блок $E_8/2401$ имеет ширину не меньше
$\sqrt{7/6}$ --- это следует из \emph{оценки через проекцию на плоскость}
(предложение~\ref{prop:planar}: для любой эйзенштейновой решётки
$D((3+\omega)\Lambda)\ge\sqrt{7/3}\,\lambda_1$) и известного радиуса покрытия
$E_8$, --- и оставляет бюджет $1/7$; его хватает на плоский блок из
$19$ цветов ($6/7+4/31=214/217<1$, откуда $45619$ в~$\R^{10}$ --- первая
оценка в $\R^{10}$ ниже классической башни~$3^n$). Ни одного численного шага
в этом доказательстве нет (\S\ref{sec:product}). Тот же бюджет с одномерным
блоком из четырёх цветов даёт $\chi(\R^9)\le9604$ при
$\ell<\sqrt{63/61}$ (следствие~\ref{cor:dim9}); теперь эта оценка вытеснена
ламинированием: $7203$ цветов --- и при этом более широкий отрезок.

Все утверждения настоящей статьи --- теоремы или точные сертификаты.
Численные конструкции с меньшим числом цветов, лестницы ширин, эйзенштейново
тождество в полной форме, строгие экраны индекса и мозаичные раскраски вынесены
в сопутствующую статью~\cite{Widths}; поисковые алгоритмы, которыми конструкции
были найдены, отрицательные экраны и журналы кампаний --- в полную
версию~\cite{Full}. Раздел~\ref{sec:repro} содержит таблицу пяти утверждений с
входными данными, верификаторами и ожидаемым выводом.
